\section{Limites du travail}
\label{sec:limites}

Cette section est la seule qui rassemble les limites du travail ; les autres chapitres y
renvoient sans les reprendre.

\paragraph{Données et généralisation.}
CICIDS2017 repose sur des scénarios construits et sur des caractéristiques de flux déjà
extraites ; les CSV d'origine sont employés, sans les corrections publiées
\cite{engelen2021troubleshooting,lanvin2023errors}. Le score dépend fortement du
protocole : le même Random Forest passe de $0{,}9951$ en séparation aléatoire à $0{,}4374$
en séparation temporelle. Dans la campagne d'origine du holdout par scénario,
l'échantillonnage retenait les premières lignes du premier fichier lu, si bien que
l'entraînement ne couvrait qu'un seul scénario d'attaque ; cette campagne a été corrigée et
réexécutée, et ses anciens chiffres ne sont plus utilisés. Après correction, aucun des cinq
modèles évalués ne dépasse la prédiction systématiquement positive
(section~\ref{sec:cicids-resultats}). La provenance des deux fichiers CSE-CIC-IDS2018 a été
vérifiée par un nouveau téléchargement depuis la source : leurs empreintes coïncident avec
les copies utilisées par les campagnes. Seules les journées des 15 et 16 février sont
cependant évaluées, et les fichiers s'arrêtent à une limite de lignes compatible avec une
exportation par tableur. Le score du 15 vers le 16 repose sur le port de destination et ne
se transfère pas dans le sens inverse (section~\ref{sec:csecic-resultats}). Le corpus
Linux/auth est un jeu public entièrement synthétique, généré avec la bibliothèque Faker et
diffusé sur Kaggle ; ses variables \texttt{status} et \texttt{service} peuvent servir de
mandataires de l'étiquette. Aucun résultat obtenu sur Linux/auth ne vaut pour des journaux
d'authentification réels. Aucun de ces résultats ne se transpose sans nouvelle évaluation à
d'autres capteurs, topologies ou périodes d'observation.

\paragraph{Corpus et incertitude.}
Le test HDFS couvre les $115\,013$ blocs de la partition finale, dont $1\,679$ positifs.
L'intervalle bootstrap du F1, $[0{,}8568\,;\,0{,}8795]$, décrit l'incertitude sur ce corpus
d'un seul système, non sur une population de systèmes distribués. Les cinq graines ne
modifient que l'échantillon d'ajustement du parseur et donnent la même matrice de
confusion : elles constituent une seule observation du test. La sensibilité du résultat à la
taille de cet échantillon, $6\,000$ blocs, n'a pas été évaluée, et la prévalence plus faible
du test ($1{,}46\,\%$ contre $3{,}48\,\%$ en validation) abaisse mécaniquement la précision
au seuil choisi en validation. Sur BGL, le groupe des templates connus ne contient aucune
anomalie et $89{,}81\,\%$ du test relève d'un template inconnu de l'entraînement ;
Histogram y applique exactement la règle « sévérité absente de l'entraînement ». Les
variantes de protocole BGL, dont l'une a été conçue après lecture du test, sont
exploratoires et ne remplacent pas le protocole principal.

\paragraph{Routage et normalisation.}
La politique open-set rejette $20/20$ fichiers inconnus mais écarte aussi $3/46$ fichiers
connus. Ses scores reprennent pour les journaux texte les colonnes \texttt{dataset} et
\texttt{subtype} du parseur de format : le rejet mesure donc en partie la capacité du
parseur à ne pas reconnaître un format. Sa calibration ne repose que sur trois fichiers
inconnus, les fichiers d'un même schéma ne sont pas indépendants et la
\texttt{decision\_margin} n'est pas une probabilité calibrée. La comparaison à une chaîne
sans routage n'a pas été exécutée ; H3 reste non tranchée. La couverture multi-format mesure
la lecture, le parsing et l'émission des unités présentes ; la voie Apache repose sur une
seule fixture, et les flux réseau tabulaires ne renseignent aucun des sept champs standard
mesurés.

\paragraph{Grandeurs non mesurées.}
Les sources ne partagent pas de vérité terrain homogène : aucune exactitude globale n'est
calculée. Aucune expérience n'analyse conjointement un même incident à partir de journaux
système et de flux réseau ; l'hétérogénéité étudiée est celle des formats, lus séparément.
La corrélation est évaluée en circuit fermé, sur des corpus synthétiques construits avec les
clés mêmes du corrélateur. L'utilisabilité du tableau de bord, le temps d'investigation et
la charge cognitive n'ont fait l'objet d'aucune étude auprès d'analystes. La priorisation
des anomalies, les traitements périodiques et l'exécution par micro-lots ne sont pas
évalués. Le budget local de temps et de mémoire associé à H4 n'a jamais été fixé, et il ne
l'est pas après coup.

\paragraph{Distribution et exploitation.}
La validation distribuée porte sur deux machines virtuelles et un Redis central. La
campagne d'endurance offre environ $0{,}5$ tâche par seconde, loin de la capacité, et la
répartition de ses tâches suit les profils de confiance déclarés. En saturation, le coude se
situe entre 8 et 16 tâches/s, avec trois expirations de délai sur $11\,430$ tâches, et la
seconde machine ne traite que $1{,}23\,\%$ des tâches réussies : la distribution est
fonctionnelle, mais elle n'accroît pas la capacité observée. Dans la détection HDFS
répartie, les utilités égales sont départagées par un ordre fixe et toutes les tâches vont
au même agent ; l'inférence n'y est pas recalculée sur les machines virtuelles. La trace de
messages de l'endurance est tronquée à $10\,001$ événements, ce qui rend ses compteurs non
conclusifs. Ni panne du broker, ni partition réseau, ni fonctionnement multi-site, ni
vieillissement sur plusieurs jours ne sont testés. TLS, contrôle d'accès, gestion des
secrets, journalisation protégée, réplication, sauvegarde et observabilité distribuée ne
sont pas évalués.

\paragraph{Mécanismes adaptatifs.}
La fonction d'utilité est un classement à poids fixes, non appris
(section~\ref{sec:spec-utilite}). La mémoire n'apporte de gain qu'à une probabilité d'échec
de $0{,}9$, avec une équité moindre, dans un protocole où la confiance déclarée est
volontairement éloignée de la fiabilité réelle (section~\ref{sec:e-h5}). Contract Net
n'obtient aucun avantage sur un répartiteur au moindre chargé
(section~\ref{sec:e-cnp-moindre}) ; ce répartiteur lit l'état des agents dans un seul
processus, sans coût de communication ni information périmée, et la comparaison n'a pas été
reprise en exécution distribuée. Sous Contract Net, quelques tâches sont perdues, parce que
la liste de repli est figée au moment de l'appel d'offres. Dans la panne appariée, l'absence
de doublon provient du magasin d'idempotence, qui l'évite aussi lorsqu'il est ajouté aux
architectures de référence. La mise à jour contrôlée des modèles dispose d'un mode
\textit{dry-run}, mais d'aucun cycle complet et répété de promotion et de rejet
(tableau~\ref{tab:statut-model-update}).

\paragraph{Statut du système.}
Ariel \logminer{} est un prototype de laboratoire. Il ne constitue ni une plateforme
opérationnelle de SOC ni un système autonome prêt pour la production.
